\documentclass[11pt]{article}
\usepackage[margin=1in]{geometry}
\usepackage{amsmath,amssymb,amsthm,mathtools}
\usepackage{enumitem}
\usepackage{array,longtable,booktabs}
\usepackage[hidelinks]{hyperref}
\usepackage{microtype}
\usepackage{xcolor}

\newcommand{\ALCIRCC}{\mathcal{ALCI}_{\mathrm{RCC5}}}
\newcommand{\EQ}{\mathsf{EQ}}
\newcommand{\PP}{\mathsf{PP}}
\newcommand{\PPI}{\mathsf{PPI}}
\newcommand{\PO}{\mathsf{PO}}
\newcommand{\DR}{\mathsf{DR}}
\newcommand{\Rel}{\mathsf{Rel}}
\newcommand{\Hor}{\mathsf{Hor}}
\newcommand{\Comp}{\operatorname{Comp}}
\newcommand{\tp}{\operatorname{tp}}
\newcommand{\code}[1]{\texttt{#1}}
\newcommand{\severity}[1]{\textbf{#1}}

\title{Referee Report on Round 13: Targets G2 and G3}
\author{Adversarial review of the D-discipline extraction and coverage walk}
\date{July 8, 2026}

\begin{document}
\maketitle

\begin{center}
\fbox{\begin{minipage}{0.92\linewidth}
\textbf{Overall verdict.}  Round 13 does not currently prove Theorem B(a) in
D-form.  The literal D2 propagation/verticalization clause is too strong for
finite-width extraction: a fixed satisfiable concept forces infinitely many
obligation-carrying vertical sources to be inserted into the same lower bag.
The Shadow Amalgamation Lemma itself remains strongly supported under the
horizontal-row discipline; the defect is on the completeness/extraction side
that round 13 explicitly designated as G2.  Status remains: strongly
supported, not certified.
\end{minipage}}
\end{center}

\section{Executive verdict table}

\begin{longtable}{p{0.11\linewidth}p{0.16\linewidth}p{0.63\linewidth}}
\toprule
Target & Verdict & Finding \\
\midrule
G2a & \severity{Fatal} & Literal D2 fallback plus D4-style liveness is incompatible with any finite
width bound.  The witness concept $C_{\mathrm{G2a}}$ below is satisfiable, but
for every prefix length $m$ the D2 extraction of the standard vertical support
chain forces the root edge bag to contain $m+1$ distinct manifestations. \\
G2b & \severity{Fatal, derivative} & The rows can be kept horizontal only by applying D2 fallback at every downward
crossing of every remote source in the tower.  Without fallback, D1 fails; with
fallback, width fails. \\
G2c & \severity{Major} & The ``no third category'' assertion is true only under the same globally
omnidirectional D2 that causes G2a.  Any width-preserving repair must replace
Coverage by a relation-indexed lemma; the present proof cannot simply be
reused. \\
G3a & \severity{Major} & Under literal D2 the coverage walk is essentially trivial but
unextractable.  Under the needed weakening, the one-sentence dismissal of
mixed walks is not a proof and lacks the invariant saying where propagation is
allowed to stop and later restart. \\
G3b & \severity{Major, conditional} & A closed three-point RCC5 network has $x\PP z$, $z\PPI y$, and $x\PP y$.
Thus a vertical endpoint relation can be supported by a mixed $\PP/\PPI$ path,
and D3 does not propagate an $\forall\PP.D$ obligation to $y$.  Literal D2
avoids this only by co-bagging $x$ with $y$, which is exactly the source of
G2a. \\
G3c & \severity{Sound under literal D2, unusable for extraction} & I did not find an independent defect in Theorem 4.3's row-realization
argument under D1--D5.  The problem is that the D2 hypothesis needed to exclude
default pairs involving obligation sources is too expensive to extract at
bounded width. \\
\bottomrule
\end{longtable}

\section{G2a counterexample: verticalization versus width}

The smallest literal-wording witness I found is
\[
C_{\min}=\exists\PP.\top\ \sqcap\
        \forall\PP.(\exists\PP.\top\ \sqcap\ \forall\PP.\exists\PP.\top).
\]
It already breaks round 13 if D2 is applied to every nonempty universal
obligation set, including purely vertical universal obligations.

To remove the likely defense ``vertical universals are meant to ride D3 and
not cast shadows'', I used the following robust version in the machine check:
\[
C_{\mathrm{G2a}}=\exists\PP.\top\ \sqcap\
\forall\PP.\bigl(\exists\PP.\top\ \sqcap\
        (\forall\PP.\exists\PP.\top\ \sqcap\ \forall\DR.A)\bigr).
\]
The added $\forall\DR.A$ makes every strict $\PP$-successor of the root carry
a genuinely horizontal universal obligation.

\subsection{Semantic forcing}

Let $a_0\models C_{\mathrm{G2a}}$.  Choose $a_1$ with $a_0\PP a_1$.  Since
$a_0$ satisfies the outer $\forall\PP$, the element $a_1$ satisfies
$\exists\PP.\top$, $\forall\PP.\exists\PP.\top$, and $\forall\DR.A$.  Choose
$a_2$ with $a_1\PP a_2$.  By $\Comp(\PP,\PP)=\{\PP\}$, $a_0\PP a_2$, so the
outer universal at $a_0$ applies to $a_2$ as well.  Inductively every finite
prefix
\[
  a_0\PP a_1\PP a_2\PP\cdots\PP a_m
\]
is forced, and every $a_i$ with $i\ge 1$ has a nonempty horizontal obligation
set because $\forall\DR.A\in\tp(a_i)$.

This is not an artefact of a chosen model.  A finite cycle cannot close the
chain: repeated $\PP$ composition would imply a nontrivial $\PP$ self-edge,
contradicting strong equality/identity.

\subsection{Why literal D2 forces unbounded bags}

Use the standard edge-bag presentation of the selected $\PP$ support chain:
$B_i$ initially contains $a_i$ and $a_{i+1}$, with adjacent bags sharing
$a_i$ or $a_{i+1}$ as usual.  Consider a remote source $a_k$ with
$k\ge 2$ and propagate its obligation shadow downward toward the root bag
$B_0$.

At the crossing from $B_i$ to $B_{i-1}$, for $1\le i<k$, the continued
separator occurrence is $a_i$.  The true relation from source $a_k$ to this
separator is
\[
  a_k\ \PPI\ a_i,
\]
because $a_i\PP a_k$.  A row entry $a_k\PPI a_i$ is illegal by D1
(horizontal rows).  Therefore literal D2 says that the crossing is replaced by
verticalization: $a_k$ is inserted into the target bag as a live port.  Repeating
this over $i=k-1,k-2,\ldots,1$ inserts $a_k$ into $B_0$.  Since this holds for
every $k$, the root edge bag must contain
\[
  \{a_0,a_1,\ldots,a_m\}
\]
for the length-$m$ prefix.  The width requirement is therefore at least $m+1$,
for arbitrary $m$, for the fixed concept $C_{\mathrm{G2a}}$.

The round-13 accounting by generation reasons (g1)--(g5), including
$B_{r+1}\le c(n)(1+B_r)^2$, counts ports generated inside a bounded side
context.  It does not count the same lower bag receiving manifestations of
unboundedly many independently generated remote obligation sources.  These are
manifestation copies, not new semantic elements, but width counts
manifestations.

\subsection{Machine check}

I added \code{verification/wp18\_g2a\_verticalization\_width\_counter.py}.
It checks three finite claims:
\begin{enumerate}[leftmargin=2em]
\item $C_{\min}$ and $C_{\mathrm{G2a}}$ are satisfiable according to the
project oracle \code{cover\_tree\_tableau.check\_sat};
\item the complete $\PP$-chain prefix network on $a_0,\ldots,a_m$ is closed
under the RCC5 composition table;
\item literal D2 fallback inserts all sources $a_1,\ldots,a_m$ into the root
edge bag.
\end{enumerate}
The relevant output was:
\begin{verbatim}
C_min_literal = (ExPP.Top n AllPP.(ExPP.Top n AllPP.ExPP.Top))
oracle_sat_min = True closure_size = 12
C_G2a_robust = (ExPP.Top n AllPP.(ExPP.Top n (AllPP.ExPP.Top n AllDR.A)))
oracle_sat_robust = True closure_size = 18

finite prefix check under literal D2:
m  chain_closed  root_bag_size  expected  fallbacks
 1         True              2         2          0
 2         True              3         3          1
 3         True              4         4          3
 4         True              5         5          6
 5         True              6         6         10
 6         True              7         7         15
 7         True              8         8         21
 8         True              9         9         28
 9         True             10        10         36
10         True             11        11         45

conclusion = PASS: root width grows unboundedly with the prefix
\end{verbatim}

\section{G2b: extracted horizontal rows}

The tower above also diagnoses G2b.  At each downward boundary, the extracted
true row entry from source $a_k$ to the continued separator $a_i$ is $\PPI$.
Thus there are only two options:
\begin{enumerate}[leftmargin=2em]
\item keep the true row entry, violating D1; or
\item apply D2 fallback, inserting $a_k$ as a live port in the target bag.
\end{enumerate}
Round 13 chooses the second option, but G2a shows that the second option has
unbounded simultaneous cost.  Therefore the manuscript has not shown that
``after saturation and threshold verticalization every shadowed pair has a
horizontal true relation'' within the claimed width bound.

This is stronger than the old overlap-amalgam stress case: no exotic side
amalgam is needed.  A single forced vertical tower with a harmless horizontal
universal at each level is enough.

\section{G2c: no third category}

The current Coverage Lemma is stated for every occurrence $y$ whenever
$\Omega_x\ne\varnothing$.  Under literal D2 this is almost tautological:
$x$ is live or shadowed in every bag reachable from its starting point, so a
manifestation of $y$ is either co-bagged with $x$ or row-covered by $x$.
However, that literal reading is precisely what fails G2a.

A width-preserving repair must allow propagation to stop at some vertical
barriers.  Once this is done, the trichotomy cannot remain occurrence-wide.  It
must become relation-indexed: for a fixed formula $\forall R.D\in\tp(x)$ and a
target $y$ with $\rho_T(x,y)=R$, either the pair is co-bagged, or it is covered
by a horizontal row if $R\in\{\DR,\PO\}$, or it is covered by a uniformly
oriented vertical support chain if $R\in\{\PP,\PPI\}$.  Pairs in a vertical
barrier that are not in relation $R$ should be allowed to remain undeclared for
that particular obligation.  The round-13 text does not state this weaker,
relation-indexed lemma, and its existing proof does not establish it.

\section{G3: the coverage walk}

\subsection{G3a: mixed walks}

The proof of Lemma 6.1 says that mixed walks ``end in cases (a)/(b) at the last
horizontal manifestation's bag family.''  I could not make this sentence into a
valid induction invariant.  A mixed walk can switch from horizontal movement to
vertical fallback, and after the fallback the source is live again; literal D2
then sends it in all directions, causing the G2a blow-up.  If this propagation
is weakened to preserve width, the proof needs an explicit rule saying which
components are still covered and which are certified irrelevant for the current
relation $R$.

Thus G3a is not an independent counterexample under literal D2, but it is not a
repairable proof as written once G2a is fixed.

\subsection{G3b: orientation bookkeeping}

I added \code{verification/wp19\_g3\_mixed\_orientation\_probe.py}.  It checks
the closed three-point network
\[
  x\PP z,
  \qquad z\PPI y,
  \qquad x\PP y.
\]
The network is composition-closed because $\PP\in\Comp(\PP,\PPI)$.
For a source type containing $\forall\PP.D$, D3 propagates $D$ and
$\forall\PP.D$ from $x$ to $z$ along $x\PP z$, but it does not propagate from
$z$ to $y$ along $z\PPI y$.  The output was:
\begin{verbatim}
B0={x,z}, B1={z,y}, overlap={z}
relations: x PP z, z PPI y, completion x PP y
composition_closed = True
D3_from_AllPP_D_at_x gives D flags: {'x': False, 'z': True, 'y': False}
D3_from_AllPP_D_at_x gives AllPP.D flags: {'x': True, 'z': True, 'y': False}
endpoint_relation_x_y = PP
conclusion = PASS: mixed PP/PPI path can have vertical endpoint; uniformity needs an extra invariant
\end{verbatim}

This is not a valid D-abstraction under literal D2, because D2 would
verticalize $x$ across the first boundary and make $x,y$ co-bagged in the
second bag.  But that observation is exactly the problem: the orientation claim
in case (c) is safe only because D2 is doing expensive global co-bagging.  It
is not a consequence of the composition table or patchwork alone.

\subsection{G3c: interaction with Theorem 4.3}

I found no direct defect in the row-realization part of Theorem 4.3 under the
literal D-discipline.  D1, D5, and the extended-family patchwork proof do make
declared horizontal row values literal values of $\rho_T$.  Also, if D2 is read
literally, default $\DR$ pairs cannot involve obligation-carrying sources,
because those sources have been propagated everywhere.

The caveat is again extraction: the literal D2 hypothesis needed for this
argument is not finite-width extractable.  After a width-preserving repair, the
``default pairs involve only obligation-free sources'' statement will need to
be replaced by a relation-indexed irrelevance statement.

\section{Steps checked and found sound}

The following pieces survived this targeted review.
\begin{enumerate}[leftmargin=2em]
\item I reran WP17.  The four finite RCC5 facts used by the Shadow
Amalgamation Lemma passed; 400/400 random D-discipline tree configurations
were extended-bag closed and globally realizable; the negative control still
triggered.
\item I reran WP16 in repaired mode.  The listed referee diagnostics and the
random sweep agreed with the cover-tree tableau oracle.  This supports the
acceptance-level intuition but does not test the finite-width extraction
failure above.
\item I reran WP14.  The two-bag patchwork checks, finite tree-patching
stress test, and strong-EQ non-forcing check passed.
\item I did not find a new defect in the horizontal Shadow Amalgamation Lemma
itself.  The D1 horizontal-row restriction plus D5 pair coherence are exactly
what make the extended-bag closure proof work.
\item The D3 propagation rule is sound for uniformly oriented vertical chains:
$\Comp(\PP,\PP)=\{\PP\}$ and $\Comp(\PPI,\PPI)=\{\PPI\}$.
\end{enumerate}

\section{Recommended round-14 repair direction}

Round 14 should not try to save literal D2.  The repair should split obligation
coverage by relation and replace omni-directional port verticalization with a
bounded horizontal-frontier mechanism.

Concretely, define horizontal obligations
\[
  \Omega_x^H=\{(R,D)\in\Omega_x:R\in\{\DR,\PO\}\}
\]
and vertical obligations $\Omega_x^V$ analogously.  Vertical obligations should
be handled by D3 along certified uniformly oriented support chains, not by
long-range shadows.  Horizontal shadows should propagate only across open
horizontal frontiers.  When a crossing is known to enter a vertical cone of the
source, the automaton should record a finite vertical-barrier certificate rather
than inserting the source as a port.  If a later side context exits the cone and
creates a possible $\DR/\PO$ target for the source, the certificate must reopen
a horizontal row there.

The corresponding Coverage Lemma should be restated as:
\begin{quote}
For every $\forall R.D\in\tp(x)$ and every $y$ with $\rho_T(x,y)=R$, the pair
$(x,y)$ is either co-bagged, row-covered by a horizontal row when
$R\in\{\DR,\PO\}$, or connected by a co-bagged uniformly oriented vertical chain
when $R\in\{\PP,\PPI\}$.
\end{quote}
This is a proposed repair path, not a completed proof.  The new proof
obligation is to show that vertical-barrier certificates are locally checkable,
width-bounded, and sufficient to reopen horizontal rows exactly when a
horizontal target can occur.

\section{Final recommendation}

Reject round 13 as a complete repair of Theorem B.  Keep the Shadow
Amalgamation Lemma as a strongly supported component, but withdraw the claimed
D-form extraction theorem and Coverage Lemma until D2 is replaced and the
relation-indexed coverage proof is supplied.

\end{document}
